Пусть $K$ -- кольцо главных идеалов, $p$ — неразложимый элемент $K$, 
$ab$ делится на $p$, $a$ не делится на $p$. 
Рассмотрим $I$ -- множество таких элементов $x \in K$, 
что $ax$ делится на $p$.

\underline{а) $I$ -- идеал.}

Если $x, y \in I$, то $x + y \in I$, так как 
$ax\ \vdots\ p, ay\ \vdots\ p \Rightarrow a (x + y)\ \vdots\ p$.

$\forall k \in K \ \ \forall x \in I \ \ kx \in I$, так как
$ax\ \vdots\ p \Rightarrow a(kx)\ \vdots\ p$.

\underline{б) $1 \not \in I, p \in I, b \in I$}

$1 \not \in I$, так как $a \cdot 1 = a\ \cancel{\ \vdots\ }\ p$.

$p \in I$, так как $a \cdot p\ \vdots\ p$.

$b \in I$, так как $ab\ \vdots\ p$.

\underline{в) $I = (p)$}

Так как $K$ -- кольцо главных идеалов, то $\exists y \in K\ :\ I = (y)$.
\newline $p \in I$, следовательно, $\exists k \in K\ :\ p = ky$.
Так как $p$ -- неразложимый, то либо

1) $k$ обратим -- Тогда $p \sim y$, то есть $I = (p)$.

2) $y$ обратим -- Но тогда идеал, порождённый $y$, совпадает со всем кольцом, а значит
содержит и единицу (противоречие с пунктом б).

\underline{г) $b$ делится на $p$}

$I$ содержит элементы, которые делятся на $p$, а $b \in I$.
Следовательно, $b$ делится на $p$.

\underline{д) Любой неразложимый элемент $K$ является простым.}

Пусть $p \in K$ неразложим, 
тогда для любых $a, b \in K\ :\ ab\ \vdots\ p$ верно что либо 
$a\ \vdots\ p$, либо $b\ \vdots\ p$ (по пункту г)).

